\documentclass[12pt]{article}
\usepackage[utf8]{inputenc}
\usepackage[T1]{fontenc}
\usepackage{hyperref}
\title{Western Neo-Aramaic Language Resources for K-12 Curriculum Development}
\author{Compiled by Claude Sonnet 5.5 (claude-sonnet-5-5) for Andr\'as Kornai}
\date{}
\begin{document}
\maketitle

\section{Introduction}
Western Neo-Aramaic (ISO 639-3 \texttt{amw}) is spoken in Ma'lula, Bakh'a and Jubb'adin in Syria. Search results describe Bakh'a as largely depopulated, Jubb'adin as the most vibrant village, and the villages as affected by the civil war. Available resources are mainly scholarly.

\section{Primary Educational Resources}
No primers, textbooks or school programmes were found.

\section{Linguistic and Reference Materials}
\begin{itemize}
\item Anas Abou-Ismail, \emph{Western Neo-Aramaic: The Dialect of Jubaadin} (Cambridge Scholars, 2019), with grammar, texts and an etymological dictionary: \url{https://www.cambridgescholars.com/product/978-1-5275-3352-3}
\item DFG project on a Dictionary of Western Neo-Aramaic (W.\ Arnold; 2019--2020); publication status was not confirmed: \url{https://gepris.dfg.de/gepris/projekt/426492122}
\item Wikipedia overview: \url{https://en.wikipedia.org/wiki/Western_Neo-Aramaic}
\item Omniglot page, with links to a chart and numeral data: \url{https://omniglot.com/writing/westernneoaramaic.htm}
\end{itemize}

\section{Religious and Community Resources}
No scripture or community organisation was found in search results.

\section{Digital and Audio Resources}
\begin{itemize}
\item Archive of written records of Western Neo-Aramaic, including Jubb'adin and Bakh'a specimens and a Syriac-script version (blog): \url{https://aymennjawad.org/24069/archive-of-written-records-of-western-neo-aramaic}
\item Endangered Languages Project entries: \url{https://endangeredlanguages.com/node/39045}, \url{https://endangeredlanguages.com/node/95862}
\end{itemize}
No freely accessible online dictionary or audio corpus was found.

\section{Curriculum Development Context}
The language is not known to be taught in schools. Any curriculum would need to build from the scholarly grammar, text collections (Arnold's published volumes, mentioned in the DFG project description) and work with remaining speakers, given the displacement of the community.

\end{document}